\documentclass[11pt]{article}
\usepackage{amsmath,amssymb,amsthm}
\usepackage[margin=1in]{geometry}
\usepackage[hidelinks]{hyperref}
\newtheorem{lemma}{Lemma}
\newtheorem{theorem}{Theorem}
\newtheorem{corollary}{Corollary}
\begin{document}
\setlength{\emergencystretch}{1em}
\title{Sector separation for the lower-five quadratic-gap family}
\author{}\date{}\maketitle
For $\beta>0$ let $\mathcal L_\beta(X^j)=(\beta)_j$.
We prove Strong Factorial for every nonzero complex
$X^5(A+BX^2)$ at every positive real Gamma shape $\beta$.
A common root would lie in a sector narrower than a zero-free
sector of the lower moment polynomial.
No half-plane stability assertion for that polynomial is needed.

Put
\[
 T_{n,\gamma}(t)=\sum_{k=0}^n\binom nk(\gamma)_{2k}t^k,\quad
 B_{0,\gamma}=1,\quad
 B_{j+1,\gamma}=(z+2j+\gamma+1)B_{j,\gamma}-zB_{j,\gamma}'.
\]
In a fixed pair let $M=\gamma+2n$, $s=2(n+1)$,
$B=B_{n,\gamma}$, $C=B_{n+1,\gamma}$ and $u=z^2$.
Write $B=B_0(u)+zB_1(u)$ and $C=C_0(u)+zC_1(u)$ for their
cyclic components; the shape $\gamma$ stays fixed in both auxiliaries.
Define
\[
 A_5=(M+2)_5,\quad b=3s(M+4)^2,\quad
 k=(M+4)(M+5)(M+6),\quad c=M+2n+8,
\]
\[
 U(u)=A_5+bu+su^2,\qquad W(u)=su(k+cu).
\]

\begin{lemma}[The general shift-five condition]\label{lem:shift-general}
For $n\ge1$, $\gamma>0$, a common root $t$ of
$T_{n,\gamma}$ and $T_{n+1,\gamma+5}$ gives $u=-1/t\ne0$ with
\[
 U(u)C_i(u)+W(u)B_i(u)=0,\qquad i=0,1.                 \tag{1}
\]
\end{lemma}
\begin{proof}
Set $R=z^{2n}T_{n,\gamma}(-z^{-2})$, $E=M-zD_z$ and
$K_j=(E)_jR$.
Pascal's identity and Gamma shifting give
\[
 (\gamma)_5z^{2(n+1)}T_{n+1,\gamma+5}(-z^{-2})
                    =[u(E)_5-(E)_7]R.
\]
Here is an explicit reduction valid for indeterminate $M,n$,
with $\gamma=M-2n$.
On a vector $(r,e,k)$ in the basis $(R,K_1,K_2)$,
application of $E$ is
\[
 (-2ur'-\gamma uk,\quad
   r-e-2ue'+uk,\quad e+Mk-2uk'),                       \tag{2}
\]
where primes mean $d/du$.
This follows from the lower Euler equation
$EK_2=MK_2+uK_1-\gamma uR$ and $EK_1=K_2-K_1$.
Starting at $(1,0,0)$ and applying $E+j$ successively,
the last two coefficients of $u(E)_5-(E)_7$ are exactly
$-W,-U$. Thus the upper zero at $R=0$ says
$WK_1+UK_2=0$.

The full rising-jet concomitant in the proof of
\href{LOWER_TWO_PRIME_GAP_SFC.tex}{Theorem \texttt{thm:mixed}},
equation(3) at $d=2$, is
$B(\pm z)K_2-C(\pm z)K_1=0$ at $R=0$.
The nonzero cubic solution jet cannot have $K_1=K_2=0$
at an ordinary point $z\ne0$.
Taking the two by two determinant with the upper relation gives
$UC(\pm z)+WB(\pm z)=0$.
Characteristic-zero even--odd inversion proves (1).
This argument does not divide by $R'$, $B$ or $U$.
\end{proof}

\begin{lemma}[A general strict Pick quotient]\label{lem:pick-general}
For $n\ge2$ and $\gamma>2n-2$, $B_1$ has only negative real
simple zeros, and for every nonreal $u$,
\[
 \frac{\Im(C_1/B_1)(u)}{\Im u}>0.                      \tag{3}
\]
\end{lemma}
\begin{proof}
The Gamma-convolution and half-plane bound, equations(3)--(4) of the
\href{CUBIC_ADJOINT_STABILITY_SFC.tex}{cubic adjoint note},
give positive coefficients for $B$ and
$\Re\zeta_j\le-\rho$, $\rho=\gamma-n+2>n$, for its roots.
For real $y$,
\[
 0<\phi'(y):=\Re\frac{B'(iy)}{B(iy)}
   =\sum_j\frac{-\Re\zeta_j}
      {(-\Re\zeta_j)^2+(y-\Im\zeta_j)^2}
 \le n/\rho<1.
\]
The continuous argument $\phi(y)=\arg B(iy)$ for $y\ge0$
starts at0 and tends to $n\pi/2$.
Consequently the zeros of $yB_1(-y^2)=\Im B(iy)$ give exactly
$\deg B_1$ positive $y$ values, all simple.
They give all zeros $-a_j$ of $B_1$, with $a_j>0$;
none is a zero of $B_0$.

Partial fractions and the positive leading coefficients give
\[
 r(u)=B_0(u)/B_1(u)=d_1u+d_0-\sum_j\frac{r_j}{u+a_j},
 \qquad d_1\ge0.
\]
At $y=\sqrt{a_j}$, differentiation gives
$\phi'(y)=2a_j/r_j$, so $r_j>2a_j$.
The raising recurrence gives
$C_1=B_0+MB_1-2uB_1'$ and hence
\[
 \frac{C_1}{B_1}
 =M+r(0)+d_1u+
       \sum_j\frac{(r_j-2a_j)u}{a_j(u+a_j)}.
\]
Thus the quotient in (3) equals
$d_1+\sum_j(r_j-2a_j)/|u+a_j|^2$.
It is strictly positive: for even $n$, degree comparison gives
$d_1>0$; for odd $n\ge3$, at least one pole occurs.
\end{proof}

\begin{lemma}[A Gamma-product zero-free sector]\label{lem:moment-sector}
Let $n\ge2$, $\gamma>2n-2$, $a=\gamma/2$ and
\[
 R(a)=\frac{4(n-1)(a+1)}{(a-n+1)^2}.
\]
If $q>R(a)$, then $T_{n,\gamma}$ has no zero with
\[
 |\arg t|<\pi-2\arctan\sqrt q.                         \tag{4}
\]
For $\gamma\ge5n$, one may take $q=40/9$; for $\gamma\ge6n$,
$q=3$; for $\gamma\ge12n$, $q=1$.
\end{lemma}
\begin{proof}
Use the monic Gamma polynomial
\[
 G_{n,\alpha}(z)=\sum_{k=0}^n\binom nk(\alpha)_kz^{n-k}.
\]
Lemma \texttt{lem:gamma} of the
\href{CUBIC_ADJOINT_STABILITY_SFC.tex}{cubic adjoint note}
writes it as $\det(zI+J_\alpha)$, where $J=D+iK$ has diagonal
$\alpha+n+1-2j$, $1\le j\le n$, and successive real symmetric
off-diagonal entries of $K$ are $\sqrt{j(\alpha+n-j)}$.
For an eigenvalue $\delta$,
\[
 \Re\delta\ge\alpha-n+1,\qquad
 |\Im\delta|\le2\sqrt{(n-1)(\alpha+1)}.
\]
The second bound follows from the quadratic form of $K$ and the
largest off-diagonal weight, which occurs at $j=n-1$ here.
For $\alpha\ge a>n-1$ the squared ratio is at most $R(a)$:
its derivative with respect to $\alpha$ is negative.
Thus every negative of a root of $G_{n,a}$ or
$G_{n,a+1/2}$ has argument of absolute value less than
$\theta=\arctan\sqrt q$.

Write $G_{n,a+1/2}(z)=\prod_j(z+\delta_j)$ and let $F_a$ be
the symmetric multiaffine polarization of $G_{n,a}$.
For
$C_n(z)=\sum_k\binom nk(a)_k(a+1/2)_kz^{n-k}$,
elementary symmetric coefficients give
\[
 C_n(z)=\left(\prod_j\delta_j\right)
                      F_a(z/\delta_1,\ldots,z/\delta_n).
\]
If $|\arg z|<\pi-2\theta$, all arguments $z/\delta_j$ lie in the
open half-plane centered at angle $\arg(z)/2$, while every root
of $G_{n,a}$ lies outside it.
The half-plane Grace--Walsh--Szego theorem makes $C_n(z)\ne0$.
We use Wagner Theorem4.1, PDF page11, Section4, in the
\href{sources/wagner-multivariate-stable-polynomials-2011.pdf}{cached proof source}.
Its exact-degree and convex circular-domain hypotheses both hold.
Pairing rising-factorial factors gives
$T_{n,\gamma}(t)=(4t)^nC_n(1/(4t))$, proving (4).

The final three choices follow by decreasing $R$ in its argument
and clearing positive denominators in
\[
 \begin{split}
 40/9-R(5n/2)&=\frac{696n+304}{9(3n+2)^2}>0,\\
 3-R(3n)&=\frac{20n+7}{(2n+1)^2}>0,\\
 1-R(6n)&=\frac{n^2+30n+5}{(5n+1)^2}>0.
 \end{split}
\]
The normalization uses roots of the monic Gamma polynomial,
not of its reciprocal. Inversion preserves the angular sector.
\end{proof}

\begin{lemma}[The common-root angle ratio]\label{lem:angle-ratio}
For $n\ge2$ and $\gamma>2n-2$, every nonreal common root of
$T_{n,\gamma},T_{n+1,\gamma+5}$ has $\Re t>0$ and
\[
 \cos^2(\arg t)>Q_n(M):=
 \frac{s(M+4)F(M)}{(M+2)(M+3)(M+2n+8)^2},\quad
 F(M)=2M^2+(6n+25)M+24n+66.                           \tag{5}
\]
The positive function $Q_n(M)$ is strictly decreasing for $M>0$.
Moreover $Q_n(8n)>3/7$ and $Q_n(14n)>1/4$.
\end{lemma}
\begin{proof}
A nonreal common root gives $H_1=UC_1+WB_1=0$ by
Lemma~\ref{lem:shift-general}.
The factor $B_1$ is nonzero there.
Also $U$ cannot vanish: otherwise $W=0$, whose only roots
are the real points0 and $-k/c$.
Thus $C_1/B_1+W/U=0$ is legitimate.
For $u=x+iy$, exact expansion gives
\[
 \frac{\Im(W(u)\overline{U(u)})}{y}
  =s\{kA_5+2cA_5x+(cb-sk)(x^2+y^2)\}.
\]
The strict Pick inequality in Lemma~\ref{lem:pick-general}
forces the bracket to be negative.
Here
$D:=cb-sk=s(M+4)F(M)>0$.
Substitute $u=-1/t$ and multiply by $|t|^2$:
\[
 kA_5|t|^2-2cA_5\Re t+D<0.
\]
It implies $\Re t>0$ and
$\cos^2(\arg t)>kD/(c^2A_5)=Q_n(M)$.

For monotonicity write $F(M)=2(M+a_0)(M+b_0)$.
Its two constants are positive, and
\[
 F(-3)=6n+9>0,\qquad
 F(-(2n+8))=-4n^2-10n-6<0
\]
give $3<a_0<2n+8<b_0$.
The logarithmic derivative of $Q_n$ is
\[
 \frac1{M+4}+\frac1{M+a_0}+\frac1{M+b_0}
 -\frac1{M+2}-\frac1{M+3}-\frac2{M+2n+8}<0,
\]
by pairing the three positive terms with the negative terms.

Let $q_{\rm num}=s(M+4)F(M)$ and
$q_{\rm den}=(M+2)(M+3)(M+2n+8)^2$.
At $M=8n$,
\[
 7q_{\rm num}-3q_{\rm den}
 =8(64n^4+1492n^3+2699n^2+1634n+318)>0;
\]
at $M=14n$,
\[
 4q_{\rm num}-q_{\rm den}
 =32(98n^4+1323n^3+1390n^2+483n+54)>0.
\]
These prove the endpoint bounds.
\end{proof}

\begin{theorem}[All positive real Gamma shapes]\label{thm:gamma-sfc}
For every $\beta>0$, every nonzero complex $X^5(A+BX^2)$
has a nonzero Gamma-$\beta$ moment in every consecutive positive pair.
Equivalently $T_{n,5n+\beta}$ and $T_{n+1,5(n+1)+\beta}$
are coprime over $\mathbb C$ for every $n\ge1$.
\end{theorem}
\begin{proof}
A real common root is impossible: for even $j$ and real $t$,
\[
 T_{j,5j+\beta}(t)=\frac1{\Gamma(5j+\beta)}
 \int_0^\infty x^{5j+\beta-1}(1+tx^2)^j e^{-x}\,dx>0.
\]
The integrand is nonnegative and positive on an interval.
For $n=1$, the lower polynomial is linear with a real root,
so no common root exists.
Let $n\ge2$, $\gamma=5n+\beta$ and $M=7n+\beta$.

If $0<\beta\le n$, then $\gamma\ge5n$ and
$Q_n(M)\ge Q_n(8n)>3/7$.
In Lemma~\ref{lem:moment-sector}, take $q=40/9$ and
$\varphi_0=\pi-2\arctan(\sqrt{40}/3)$.
Here $\cos\varphi_0=31/49$, and
$3/7-(31/49)^2=68/2401>0$.
Lemma~\ref{lem:angle-ratio} would put a common root strictly inside
this zero-free sector, a contradiction.

If $n\le\beta\le7n$, then $\gamma\ge6n$,
$Q_n(M)\ge Q_n(14n)>1/4$, and $q=3$ gives a zero-free sector
$|\arg t|<\pi/3$. The angle ratio again puts a common root
strictly inside that sector.
If $\beta\ge7n$, then $\gamma\ge12n$ and $q=1$ excludes every
root with $\Re t>0$, already forced by Lemma~\ref{lem:angle-ratio}.
All positive $\beta$ are covered.

For $AB\ne0$ the exact phase normalization is
\[
 \mathcal L_\beta((X^5(A+BX^2))^j)
   =A^j(\beta)_{5j}T_{j,5j+\beta}(B/A).
\]
These factors are nonzero.
If a coefficient vanishes, the phase is a nonzero monomial and
all positive moments are nonzero.
\end{proof}

\begin{corollary}\label{cor:ordinary-sfc}
Every nonzero ordinary complex $X^5(A+BX^2)$ satisfies Strong
Factorial at every positive consecutive pair.
In particular its consecutive normalized moment polynomials are
coprime over $\mathbb Q$ and over $\mathbb C$ at every index.
\end{corollary}
\begin{proof}
Take $\beta=1$ in Theorem~\ref{thm:gamma-sfc}.
\end{proof}

The written operator, spectral, polarization and Pick arguments
supply the unbounded quantifiers. Exact controls verify formulas
and normalizations; they do not promote an index prefix.
Other supports and unrestricted FC/SFC remain separate.
Internal audit and replay are distinct from independent whole-proof
replay, formal verification, external review and priority.
\end{document}
